\chapter{Auditoría de fuentes y trazabilidad}\label{app:sources}
\section{Clases de evidencia}
La presentación y guía ULE fundamentan organización docente. NASA, JPL, MIT y ESA aportan documentación técnica pública consultada temáticamente. Las derivaciones originales se comprueban con conservación, álgebra o cálculo; la IA redactora no constituye una fuente de verdad. Una referencia descargada no equivale a un documento leído íntegramente.

En el proyecto nativo, las fuentes que no han recibido una auditoría académica de autoridad permanecen PENDING. Los fragmentos exactos registrados permiten localizar texto en una versión con hash; no certifican correspondencia semántica de todas las frases del capítulo. La declaración de suplemento abierto tampoco significa revisión científica.
\section{Mapa de cobertura}
\begin{longtable}{p{.12\linewidth}p{.37\linewidth}p{.40\linewidth}}\toprule Bloque&Capítulos editoriales&Base consultada y desarrollo\\\midrule\endhead
I&1--4: balances, termodinámica y parámetros&Guía ULE; NASA empuje e impulso; MIT 2--3. Derivación propia de balances e incertidumbre.\\
II&5--9: toberas y expansión&MIT 3 y 14; RP-1311. Área-Mach, crítico, CF, choque y comparación de modelos.\\
III&10--13: sólido&MIT 16--18 para mecanismos; balances normalizados y geometría derivados, sin trasladar formulaciones.\\
IV&14--20: líquido&MIT 9--10, 19--27 y SP-125, consulta temática; RP-1311 y documentación CEA.\\
V&21--24: eléctrico&Goebel y Katz, capítulos 1--3 y secciones de haces; conservación de energía y momento.\\
VI&25--28: vehículo&MIT 32, 33 y 35--36, guía y fichas NASA/ESA. Masa variable, etapas y trayectoria derivados con hipótesis.\\\bottomrule\end{longtable}
\section{Discrepancias y decisiones}
La página NASA de impulso específico contiene una agrupación dimensionalmente incompatible en una de las expresiones de velocidad equivalente: suma rapidez y fuerza antes de dividir por caudal. Se utiliza la relación derivada $c=v_e+(p_e-p_a)A_e/\dot m$, contrastada con el balance y con otras fuentes. El prestigio institucional no sustituye un control dimensional.

En las notas manuscritas se localizaron ambigüedades de ``mitad'' de una geometría, unidades no declaradas de una ley empírica y discrepancias entre algunos estados térmicos y velocidades calculadas. La auditoría separada distingue errores confirmados de lecturas o referencias insuficientes. Ninguna solución estudiantil se incorporó como patrón verificado.

Las notas de expansión reactiva explican la comparación congelada-equilibrio bajo pérdidas despreciadas. El manual restringe ese razonamiento a sus hipótesis: no presenta ambos ideales como cotas universales de prestaciones reales con calor, fricción o separación.
\section{Registro consultable}
El paquete editable contiene inventarios con autor, título, URL, año, hash, localizador y estado. También conserva el análisis de competencias, la matriz de ecuaciones y los informes de Studium. No se incluyen escaneos ni el PDF de Wuolah, ni sus extracciones extensas. Los originales privados permanecen en la carpeta de trabajo del usuario y no forman parte de la distribución.
